%--------------------------------------------------------------------
\chapter*{Notazioni impiegate}
\addcontentsline{toc}{chapter}{Notazioni impiegate}
\setlength{\parindent}{2pt}

\begin{multicols*}{2}
      Impiegheremo caratteri latini ($x$, $y$, $z$) per indicare quantità reali;
      caratteri latini sottolineati ($\vec{x}$, $\vec{y}$, $\vec{z}$) per indicare
      vettori o campi vettoriali; caratteri greci minuscoli ($\alpha$, $\beta$, $\gamma$) per indicare
      curve; e caratteri greci maiuscoli ($\Sigma$) per indicare superfici.

      \section*{Analisi matematica}
      \addcontentsline{toc}{section}{Analisi matematica}

      \begin{itemize}
            \item $f_i$ -- nel caso di una funzione $f : \RR^n \to \RR^m$, la proiezione di $f$ sulla $i$-esima coordinata, ovverosia
                  $\pi_i \circ f : \RR^n \to \RR$.
            \item $\der{f}{t}(x)$, $f'(x)$ -- derivata di una funzione $f : I \subseteq \RR \to \RR^n$. Nel caso $n > 1$, coincide con
                  il vettore $({f_i}'(t))_i)$. La notazione può essere iterata per ottenere le derivate successive.
            \item $\partial_{x_i} f(\vec{x})$, $\pd{f}{x_i}(\vec{x})$, $f_{x_i}(\vec{x})$ -- derivata parziale nella $i$-esima coordinata
                  di una funzione $f : \RR^n \to \RR$ nel punto $\vec{x}$. La notazione può essere iterata per ottenere le
                  derivate successive.
            \item $\vec{x_i}$ -- derivate parziali nella $i$-esima coordinata di un campo $\vec{x}$ da $\RR^n$ a $\RR^m$, ovverosia
                  $((x_1)_i, \ldots, (x_m)_i)^\top$.
            \item $\nabla f(\vec{x})$ -- gradiente di una funzione $f : \RR^n \to \RR$, ovverosia il vettore $(\partial_{x_i} f(\vec{x}))_i^\top$.
            \item $J f(\vec{x})$ -- jacobiano di una funzione $f : \RR^n \to \RR^m$ nel punto $x$, ovverosia la matrice
                  $(\partial_{x_j} f_i (\vec{x}))_{i, j} = (\nabla f_i(\vec{x}))_{i}$.
            \item $J_{\vec{y}} f(\vec{p})$ -- nel caso di una funzione $f : \RR^{m} \times \RR^{n} \to \RR^{n}$,
                  la sottomatrice $n \times n$ quadrata di $J f(\vec{p})$ date dalle ultime $n$ colonne. Coincide
                  con $J (\pi_{\RR^n} \circ f)(\vec{p})$.
            \item $C^n$ -- classe delle funzioni dotate di derivate parziali continue fino all'ordine $n$. Per $n = 0$, coincide con la classe delle funzioni continue ($C^0$).
            \item $C^\infty$, liscio -- classe delle funzioni derivabili parzialmente per un numero arbitrario di volte con continuità.
            \item diffeomorfismo di classe $C^k$ -- funzione di classe $C^k$ con inversa di classe $C^k$.
      \end{itemize}

      \section*{Geometria}
      \addcontentsline{toc}{section}{Geometria}

      \begin{itemize}
            \item $B_R(P, \RR^n)$, $B_R(P)$ -- la palla $n$-dimensionale di raggio $R$ e centro $P$. Ometteremo $\RR^n$ quando la dimensione si deduce dal contesto.
            \item $S_a^i(P)$ -- l'ipersfera $i$-dimensionale di raggio $a$ e centro $P$, ovverosia $\{\vec{x} \in \RR^{i+1} \mid \norm{\vec{x} - p} = p\}$.
            \item $\TT_{a, b}$ -- toro di raggio maggiore $a$ e raggio minore $b$.
      \end{itemize}

\end{multicols*}
